\chapter{System Design}
\label{chap:design}

\section{Introduction}

This chapter presents the architecture of CANvisionNative: the technology choices, the overall layered architecture, the desktop client, the backend, the database, the \gls{rest} \gls{api}, the replay engine, the detection pipeline, the \gls{ai} explanation engine, the hardware layer, and the deployment architecture. Every diagram in this chapter is a UML diagram produced with PlantUML, generated directly from the actual architecture described in Chapter~\ref{chap:impl}.

\section{Technologies used}
\label{sec:tech-used}

Figure~\ref{fig:tech-logos} shows the full technology stack, as presented during the project defense. Each technology is introduced individually below: what it is, and why it was the right fit for this platform.

\begin{figure}[H]
\centering
\includegraphics[width=0.85\textwidth]{tech_logos.png}
\caption{Technology stack used across the desktop client, backend, machine learning, and AI explanation layers.}
\label{fig:tech-logos}
\end{figure}

\textbf{C\# and WPF (.NET Framework 4.8).} WPF is Microsoft's UI framework for native Windows desktop applications, paired here with the \gls{mvvm} pattern. It was chosen because it gives a responsive, native Windows interface with strong data-binding support, well suited to a chart-heavy, multi-screen client, and .NET 4.8 was already the required target platform for this project.

\textbf{Python and FastAPI.} FastAPI is a Python web framework for building typed, self-documenting \gls{rest} APIs. It was chosen because it integrates directly with the Python machine-learning stack (scikit-learn, pandas) already used for model training, and it exposes a clean \gls{api} with minimal boilerplate.

\textbf{SQLite.} SQLite is a lightweight, embedded database that runs in-process, with no separate server to install or manage. It was chosen because the system runs as a single local process on a single workstation, so a networked database server would add an operational dependency with no real benefit.

\textbf{scikit-learn.} scikit-learn is a Python machine-learning library providing classical, well-supported algorithms. It was chosen because the platform's anomaly-scoring and state-classification needs (Isolation Forest~\cite{liu2008isolation}, K-Means~\cite{macqueen1967methods}) are well served by these classical methods, without the added weight of a deep-learning framework.

\textbf{asammdf and python-can.} \texttt{asammdf}~\cite{asam_mdf} decodes ASAM \gls{mf4} binary log files, and \texttt{python-can} provides a common interface to physical \gls{can} adapters. Both were chosen as the standard, actively maintained open-source libraries for these two data-access needs in the Python ecosystem.

\textbf{Google Gemini and Ollama.} Gemini is a cloud-hosted large language model, and Ollama runs a language model locally. The platform uses Gemini as its primary explanation provider, with a local Ollama model as an offline fallback, so an analyst can still get a natural-language explanation without an internet connection.

\textbf{PDFsharp.} PDFsharp is a .NET library for generating PDF documents. It was chosen for the export feature because it lets the client build a formatted report directly from C\#, without depending on an external PDF service.

\textbf{Communication.} The desktop client and the backend communicate over \gls{http} \gls{json} on the local machine, which keeps the two independently testable and lets the backend be driven directly for automated validation, as done throughout Chapter~\ref{chap:validation}.

\section{Overall architecture}

CANvisionNative is built as two cooperating processes on the same machine: a Windows desktop client and a local Python backend, communicating over \gls{http} on \texttt{localhost}. Figure~\ref{fig:overall-arch} shows the layered view of the system.

\begin{figure}[H]
\centering
\includegraphics[width=0.65\textwidth]{architecture_overall.png}
\caption{Overall layered architecture of CANvisionNative.}
\label{fig:overall-arch}
\end{figure}

The desktop client only talks to the \gls{rest} \gls{api} layer. The \gls{api} layer routes requests to the detection engine and to the data layer, and the detection engine calls the \gls{ai} explanation engine when an alert explanation is requested. This separation means the detection logic can be exercised directly through the \gls{api}, independently of the desktop client, which is how a large part of the validation in Chapter~\ref{chap:validation} was carried out.

\subsection{Frontend architecture}

The desktop client follows the \gls{mvvm} pattern. Each screen (Home, Dashboard, Telemetry, Diagnostics, Log Playback, Anomaly Intel, Settings) has a XAML view and a corresponding ViewModel.

Two services sit between the ViewModels and the backend, as shown in Figure~\ref{fig:frontend-arch}. An \gls{api} client service is a thin \gls{http} client wrapping every backend endpoint. A shared data service holds the in-memory client state (current snapshot, alert history, replay runtime metrics); it owns a polling timer that periodically calls the backend and raises events the ViewModels subscribe to.

\begin{figure}[H]
\centering
\includegraphics[width=0.75\textwidth]{architecture_frontend.png}
\caption{Frontend architecture: MVVM views bound to ViewModels, backed by two shared services.}
\label{fig:frontend-arch}
\end{figure}

Navigation between sections is driven by a section catalog and a navigation service, so that adding a new screen does not require changes to existing screens.

\subsection{Backend architecture}

The backend is a single FastAPI application organized around five responsibilities:

\begin{enumerate}
    \item \textbf{Decoding} --- vehicle profile registry and \gls{dbc}-based signal decoding.
    \item \textbf{Data processing} --- log parsing and \gls{mf4}-to-\gls{csv} conversion.
    \item \textbf{Machine learning} --- feature engineering, the runtime detection engine, training scripts, and the \gls{ai} explainer.
    \item \textbf{Hardware} --- the live \gls{can} adapter interface and session recorder.
    \item \textbf{Persistence} --- the SQLite session and alert store.
\end{enumerate}

Long-running work, such as a replay or an offline analysis, runs in a background subprocess or thread so the \gls{api} stays responsive; the client polls a status endpoint until the work completes.

\subsection{Database}

The backend keeps a small SQLite database for session bookkeeping and historical alerts. SQLite was chosen because the whole system runs as a single local process on a single workstation, so a networked database server would only add an operational dependency with no real benefit. The database stores three things: each recorded session (a live, replay, or injection run), every captured frame within that session with its normal/anomalous classification, and the history of raised alerts. This is enough to support session bookkeeping, historical alert review, and re-training from analyst-corrected labels, without needing anything beyond a local, file-based store.

\section{REST API}

The backend exposes 49 endpoints under a small number of prefixes: \texttt{/api/replay/*} (start, stop, status, alerts), \texttt{/api/live/*} (simulation and hardware capture), \texttt{/api/explain/\{index\}} (\gls{ai} explanation), \texttt{/analyze} (offline batch analysis), \texttt{/api/vehicles} (profile registry), \texttt{/api/session/*} (recording), \texttt{/api/db/*} (database access), and \texttt{/api/settings}. The full endpoint list is given in Appendix~\ref{app:api}. All endpoints return \gls{json} and are consumed exclusively by the desktop client over \texttt{localhost}.

Figure~\ref{fig:rest-api-flow} groups these endpoints by the backend component they ultimately reach: the replay, live, and offline-analysis endpoints all drive the detection engine, the explanation endpoint drives the \gls{ai} explanation engine, and the vehicle, session, and settings endpoints read and write the data layer directly.

\begin{figure}[H]
\centering
\includegraphics[width=0.75\textwidth]{rest_api_flow.png}
\caption{REST API flow: client requests grouped by the backend component they reach.}
\label{fig:rest-api-flow}
\end{figure}

\section{Replay engine}
\label{sec:design-replay}

The replay engine is launched as a subprocess by the \texttt{/api/replay/start} endpoint. It reads the input file through a single, format-aware loader (Section~\ref{sec:impl-parser}) shared with the offline batch analyzer, so that every supported format enters the pipeline as the same canonical \texttt{(timestamp, can\_id, payload)} stream. Each frame is then passed, in order, to the same \texttt{RealtimeEngine} used by the live and simulated paths, guaranteeing that replay, live, and simulated detection behave identically for the same input. Figure~\ref{fig:seq-replay} shows the full replay lifecycle as a sequence diagram. Section~\ref{sec:impl-replay} describes its implementation in detail, and Chapter~\ref{chap:validation} documents two defects that were found and fixed in this exact module.

\begin{figure}[H]
\centering
\includegraphics[width=0.9\textwidth]{sequence_replay.png}
\caption{Sequence diagram: replay lifecycle, from client request to published alerts.}
\label{fig:seq-replay}
\end{figure}

\section{Detection pipeline}
\label{sec:design-pipeline}

Figure~\ref{fig:pipeline} shows the per-frame detection pipeline, common to replay, simulation, and live hardware capture.

\begin{figure}[H]
\centering
\includegraphics[width=\textwidth]{detection_pipeline.png}
\caption{Per-frame detection pipeline: feature extraction, state classification, model selection, four-layer scoring, fusion, and alerting.}
\label{fig:pipeline}
\end{figure}

Feature extraction computes, per \gls{can} identifier, a rolling set of statistics: inter-arrival time and its variance, message frequency, byte-level difference from the previous frame on the same identifier, and Shannon entropy of the identifier stream and of the payload. Before scoring, the engine classifies the vehicle's current operating state (parking, idle, driving) with a K-Means model, shown as a state diagram in Figure~\ref{fig:state-vehicle}, and selects the Isolation Forest trained for that state and for the currently selected vehicle profile.

\begin{figure}[H]
\centering
\includegraphics[width=0.6\textwidth]{state_vehicle.png}
\caption{Vehicle driving-state classification: Parking, Idle, and Driving, classified by K-Means on 14 rolling features after RobustScaler normalization.}
\label{fig:state-vehicle}
\end{figure}

The four detection layers are then computed:

\begin{itemize}
    \item The \textbf{timing layer} flags both abnormally bursty timing and abnormally high message frequency (a pure timing-variance metric alone would miss a steady flood, since a perfectly periodic flood has almost no timing variance).
    \item The \textbf{payload layer} flags abrupt, out-of-pattern byte changes.
    \item The \textbf{ML layer} scores the feature vector with the Isolation Forest selected for the classified state and vehicle profile, falling back to a lightweight heuristic if no matching per-vehicle model is available.
    \item The \textbf{temporal-persistence layer} tracks, over a rolling window, how often the fused score has stayed above the high-severity threshold, so an isolated single-frame spike is treated differently from a sustained anomaly.
\end{itemize}

The four layers are combined into one fusion score by a fixed weighted sum, and the fused score is compared against four severity thresholds (LOW, MEDIUM, HIGH, CRITICAL). Only frames whose fused score crosses the configured alert threshold are turned into alerts, subject to a per-identifier debounce so that a single burst is not reported as dozens of near-duplicate alerts. The exact weight and threshold values are listed in Table~\ref{tab:app-fusion-config} (Appendix~\ref{app:config}).

\section{Hardware layer}

The hardware layer wraps the \texttt{python-can} library to support three transports: ELM327 over a serial/slcan connection, PEAK PCAN-USB, and native SocketCAN. It exposes connect, disconnect, and status endpoints, and feeds every received frame through the same \texttt{RealtimeEngine} instance used by replay, with the vehicle identifier selected in the Settings view. This design means the detection pipeline itself required no change to support live hardware once an adapter is available, only the transport layer differs, which is why the code-complete-but-hardware-blocked status described in Chapter~\ref{chap:requirements} was an acceptable risk to carry through the project.

\section{AI explanation engine}
\label{sec:design-ai}

Figure~\ref{fig:ai-workflow} shows the explanation workflow triggered when the analyst selects an alert.

\begin{figure}[H]
\centering
\includegraphics[width=0.75\textwidth]{ai_workflow.png}
\caption{AI explanation workflow: context building, LLM provider chain, and rule-based fallback.}
\label{fig:ai-workflow}
\end{figure}

The context builder assembles four sources before calling the language model: the current alert's own features (identifier, frequency, entropy, fusion-layer scores), the expected statistics for the vehicle's current state, a keyword match against a small attack knowledge base, and the most similar historical anomaly found by cosine similarity against a persistent, growing store of previously explained anomalies. The assembled context is sent to a language-model provider chain Gemini first, a local Ollama model second and, if neither is reachable, a rule-based template still produces a complete, correctly parameterized explanation instead of failing the request, satisfying requirement NFR6.

\section{Class diagram (backend detection core)}

Figure~\ref{fig:class-diagram} shows the simplified class diagram of the backend's detection core, the part of the system most directly involved in the validation work of Chapter~\ref{chap:validation}.

\begin{figure}[H]
\centering
\includegraphics[width=0.85\textwidth]{class_detection_core.png}
\caption{Simplified class diagram of the backend detection core.}
\label{fig:class-diagram}
\end{figure}

\texttt{RealtimeEngine} is the single entry point used by replay, live simulation, and hardware capture. It owns the current \texttt{vehicle\_id}, delegates timing/payload scoring and fusion to \texttt{ScoringEngine}, and delegates ML scoring to \texttt{RuntimeModelStore}, which holds one \texttt{StateModelBundle} per vehicle profile (plus a generic fallback bundle) and is responsible for selecting the correct per-state, per-vehicle model. \texttt{IDSStatistics} accumulates per-run counters and severity/alert-reason distributions and writes the final summary report consumed by both the desktop client and the \gls{ai} explanation engine.

\section{Deployment diagram}

Figure~\ref{fig:deployment} shows how the system is deployed on a single workstation, with the two external systems it may reach over the network.

\begin{figure}[H]
\centering
\includegraphics[width=0.85\textwidth]{deployment.png}
\caption{Deployment diagram: single-workstation deployment with external LLM and hardware links.}
\label{fig:deployment}
\end{figure}

\section{Conclusion}

This chapter presented the layered architecture of CANvisionNative, its desktop and backend components, its database schema, its \gls{rest} \gls{api}, the replay engine, the four-layer detection pipeline, the hardware transport layer, and the \gls{ai} explanation workflow, together with the class, entity-relationship, and deployment diagrams that summarize the design. The next chapter describes how this design was implemented.
